% =====================================================================
\section{The state of $H1$ after the certificate day: one strict certificate, eight strict
lemmas, one window at two grades, and twelve roads that end}
\label{ch:h1-status}
% =====================================================================

Chapter~\ref{ch:scalar-h1} reduced the family $\{H1(\delta)\}$ to the single scalar statement
$G''(0)\ge\delta$ \eqref{eq:scalar-rule} and exhibited a solvable model in which it fails. A ledger
therefore has one question left to answer: \emph{is the sign of $G''(0)$ ever \emph{known}, on any
model, without sampling?} This chapter records the day on which the answer became yes for exactly one
object, and records with equal care the five places where it is still no.

The five deliverables are these, and they do not overlap.

(i) \emph{One strict certificate} \LT{}/\LC: on a single plaquette, $\mathrm{SU}(2)$, $\beta=2.5$,
$a=0$, the gap $S-D$ is bounded away from zero by interval arithmetic at three values of the boundary
holonomy, and --- through a monotonicity that is proved in its identities and measured in its sign ---
on the whole interval $\theta\in[0,\pi]$ (\S\ref{sec:h1s-cert}). (ii) \emph{Eight strict lemmas}
\LT{}: a curvature floor $\Phi''\succeq(N/12)\mathrm{id}$ valid for every $X$ and every direction, the
sign of the mean curvature, a Stein identity that turns the variance--curvature comparison into a
square, a pointwise ratio that turns the criterion into one moment inequality, two moment inequalities
with three independent proofs each, the skeleton identity of the whole reduction, and the exact
$O(a)$ splitting that replaces an approximate affinity (\S\ref{sec:h1s-lemmas}). (iii) \emph{One
window at two grades} \LT{}/\LC: the analytic window $\beta\sin^{2}(\theta/2)\le\cos(\theta/2)$, and a
measured window that is $2.3$ times wider and whose tightness is one unproved monotonicity
(\S\ref{sec:h1s-window}). (iv) \emph{One lattice ledger} \LE: $G''(0)>0$ at every measured coupling
on $L=2$ and at $L=3$, with the sign carried by $16/16$ and $8/8$ chains (\S\ref{sec:h1s-lattice}).
(v) \emph{Twelve roads that end} \LN: twelve classical or structural routes, each with the number that
kills it (\S\ref{sec:h1s-roads}).

And one boundary, stated before any of it is read as progress (\S\ref{sec:h1s-limit}): the certificate
is a $0$-dimensional toy at one coupling with $a=0$; the physically relevant volume $L=2$ is
$143$-dimensional and has no certificate at all; the $\mathrm{SU}(3)$ gap is measured and not proved;
and $a\neq0$ has no criterion, only an exact splitting. \textbf{Nothing here proves a gap and nothing
here opens a route. Clay remains $0\%$.}

\medskip
\noindent\textbf{\LN{} \emph{What these results do not imply, stated before they are read as progress.}}
The results of this chapter establish positivity of $G''(0)$. \textbf{They do not imply a mass gap.} Two
\LT{}-level obstructions are recorded in \code{docs/paper-v5} \S2--\S3: (i)
$G''=\Delta_{\mathrm{dyn}}+W_{\mathrm{full}}$ with $W_{\mathrm{full}}\ge N/12>0$, so $G''>0$ is strictly
weaker than $\Delta_{\mathrm{dyn}}>0$; (ii) the harmonic-oscillator~$\otimes$~free-mode counterexample has
$G''=2\omega>0$ while $m=0$. The correct target is $f_{1}^{2}$ (the $O=\tfrac14\operatorname{tr}F^{2}$
pole residue) together with $(CP)$, see \code{paper-v5} \S7.

\src{\code{docs/paper-v5} \S2--\S3, whose second obstruction is M6 (``harmonic oscillator
$\otimes$ free mode'': $G''=2\omega>0$ while $m=0$), and \S7 (the target $f_{1}^{2}$ together with
$(CP)$); \code{T01-IMPLY.md}, \code{T03-TOY.md}. Nothing below is retracted: every label of this chapter
stands as stated. What is withdrawn is one \emph{reading} of them --- that positivity of $G''(0)$ is
itself a route to $m>0$.}

\subsection{The certificate: $S-D>0$ on a single plaquette}
\label{sec:h1s-cert}

\subsubsection*{The object, and the two exact simplifications that make it finite-dimensional}

Take one plaquette, $U=\exp(\xi\cdot\sigma/2)$ with $|\xi|<2\pi$, boundary holonomy
$B=\exp(i\theta\sigma_{3}/2)$, $v=e_{3}$, and condition on the slice $\xi_{3}=a=0$; write
$u=\rho/2$, $c=\cos(\theta/2)$, $s=\sin(\theta/2)$ and
$g(u)=\sin u/(2u)$. Then
\begin{equation}
G=\partial_{\xi_{3}}S_W=\beta s\,g(u),\qquad
H=\partial^{2}_{\xi_{3}}S_W=\tfrac{\beta c}{2}\,g(u),\qquad
\Phi=-\log J,\quad J=\frac{\sin^{2}u}{u^{2}} ,
\label{eq:h1s-GH}
\end{equation}
and two facts are exact for \emph{every} $\theta$, not merely at $\theta=\pi$:
\begin{enumerate}[label=(\roman*)]
\item $\Phi$ depends on $\rho$ alone and the slice $\xi_{3}=0$ is tangent to the sphere of constant
  $\rho$; hence $\Phi'=\Phi_{\rho}\,\xi_{3}/\rho\equiv0$ pointwise, so the conditioned variance has
  \emph{no} cross term and $\operatorname{Var}(G+\Phi')=\operatorname{Var}(G)$;
\item the weight on the slice collapses to one dimension,
  $\nu\propto e^{-\beta(1-c\cos u)}\sin^{2}u\,u^{-1}\,\mathrm{d} u$ on $(0,\pi)$.
\end{enumerate}
Consequently
\begin{equation}
S:=E[H+\Phi'']=\tfrac{\beta c}{2}\,E_{\nu}[g]+E_{\nu}[\Phi''],\qquad
D:=\operatorname{Var}(G)=\beta^{2}s^{2}\operatorname{Var}_{\nu}(g),
\qquad
\Phi''(u)=\frac{1}{2u^{2}}-\frac{\cot u}{2u} ,
\label{eq:h1s-SD}
\end{equation}
with the gap of the whole reduction being $\mathrm{gap}(\theta):=S-D=G''(0)$ \eqref{eq:scalar-identity}.
The first of (i)--(ii) is the reason the certificate is cheap: at $\theta=\pi$ the action on the slice
is the constant $\beta$, the measure is \emph{exactly} Haar, $H\equiv0$ pointwise, and
$E[\Phi'']=1/(4Z)=1/(2\,\mathrm{Cin}\,2\pi)=0.20511530$ is an integration by parts away from exact
($\Phi''w=-\tfrac14\mathrm{d}[\sin^{2}u/u^{2}]/\mathrm{d} u$). \LT

\src{\code{research/v3/species/Q11-CERT.md} \S1.1--1.4 (the model, the three derivatives,
$\Phi'\equiv0$, the four entire integrands, and the exact $NS=1/4$); \code{R03-CERT-THETA.md} \S1 (the
same reduction for general $\theta$, and the explicit correction that $\Phi'\equiv0$ holds for all
$\theta$); \code{Q04-WEYL.md} (2.1) and \S1.3 (that the flat-chart measure of the model \emph{is} Haar,
$\mathrm{d} U=(1/4\pi^{2})J\,\mathrm{d}^{3}\xi$ exactly, so Weyl's formula applies rather than being analogous);
\code{ch18-scalar-h1.tex} \eqref{eq:scalar-identity} for the identification $\mathrm{gap}=G''(0)$.}

\begin{claimbox}{cT}
\claimhead{Theorem}{(three certified values, $S-D>0$)}
\emph{On a single plaquette, $\mathrm{SU}(2)$, $\beta=2.5$, slice $\xi_{3}=0$, $v=e_{3}$,
$\|v\|=1$, the interval arithmetic described below certifies}
\begin{equation}
\mathrm{gap}(0)=0.729924128620548449,\qquad
\mathrm{gap}(\pi/2)=0.5429439223536897638,\qquad
\mathrm{gap}(\pi)=0.1359947110419882110,
\label{eq:h1s-points}
\end{equation}
\emph{each with a rigorous enclosure of width $\pm10^{-55}$, and in each case the lower end of $S$
exceeds the upper end of $D$; at $\theta=0$ one has $D\equiv0$ exactly, because $s=0$ makes $G$
vanish pointwise.} \LT{}
\end{claimbox}

\src{\code{R10-AUDITCERT.md} \S0 and \S2 (the three values, the enclosures, the three independent
numerical roads: its own $256\times20$ Gauss--Legendre with Cauchy remainder
$\le3.1\times10^{-82}$, \code{mp.quad} at $\mathrm{dps}=45$ agreeing to $10^{-30}$, and
\code{q04\_weyl.moments} agreeing to $10^{-16}$); \code{Q11-CERT.md} \S2.2 and \S5 (the
$\theta=\pi$ value, the intervals, and $\beta^{*}$); \code{R03-CERT-THETA.md} \S2.2 and \S3.1 (the
first certificate at $\theta\neq\pi$); \code{R06-GAPMONO.md} \S4 (the same three values used as an
independent order check on the monotonicity).}

\begin{center}
\footnotesize
\begin{tabular}{@{}lrrr@{}}
\toprule
$\theta$ & $S$ & $D$ & $\mathrm{gap}=S-D$\\
\midrule
$0$      & $0.729924128620548449$ & $0$ (exact) & $\mathbf{+0.729924128620548449}$\\
$\pi/2$  & $0.5581836427245872055$ & $0.0152397203708974417$ & $\mathbf{+0.5429439223536897638}$\\
$\pi$    & $0.2051152971230731480$ & $0.0691205860810849369$ & $\mathbf{+0.1359947110419882110}$\\
\bottomrule
\end{tabular}
\end{center}

\subsubsection*{What ``strict'' means here, and the two bugs the standard caught}

The pipeline has five ingredients and each one is machine-checked rather than asserted
(\code{Q11-CERT.md} \S2.1, \code{R03-CERT-THETA.md} \S2.1): the Taylor coefficients of the four
entire integrands are \emph{exact rationals} (\code{Fraction}; at $\theta=\pi/2$ in $\mathbb{Q}[\sqrt2]$,
never floating point); the integration is term by term,
$\int_{0}^{\pi}u^{k}\,\mathrm{d} u=\pi^{k+1}/(k+1)$, which is an identity for entire functions and not a
quadrature; $\pi$ and $\sqrt2$ enter only as $800$-bit \code{libmp} intervals with directed rounding
(\code{'f'} down, \code{'c'} up); every arithmetic operation carries an outward padding
($2^{-400}$ absolute in \code{q11}, relative in \code{r03}); and the truncation is bounded by a Cauchy
estimate, $K=260$ and $K=400$ respectively. The one input that is not pure arithmetic is
\code{libmp}'s rounding behaviour, and the padding covers it by more than a hundred orders of
magnitude.

The evidence that the directed rounding is \emph{actually} happening is the width pattern, not the
endpoints: \code{q11}'s width for $Z$ is $2.018\times10^{-118}$, and
$2.018\times10^{-118}/(2^{-400})=522=2(K_{\max}+1)$, i.e.\ exactly the per-operation accumulation of the
pad. A "high precision, roughly" computation would give width $0$ (identical endpoints) or ulp noise,
never that number. Two genuine bugs were found and fixed by using $\theta=\pi$ as a regression: the
power index of $\sin^{4}u/(4u^{3})$, which corrupted only $\mathrm{Var}(g)$ and was caught because the
regression reproduces \code{q11}'s $0.0110592937$; and the \emph{absolute} padding, which under the
catastrophic cancellation of the $\theta\neq\pi$ series (peak terms $\sim10^{7}$, coefficients as small
as $8.5\times10^{-227}$) inflated $Z$ to $\pm1.56\times10^{76}$ with the endpoints inverted --- the fix
is a purely relative pad, after which the $\theta=\pi/2$ width ($3.5\times10^{-21}$) matches the
magnitude of its own Cauchy tail ($1.77\times10^{-21}$) instead of collapsing to a point. \LT

\src{\code{Q11-CERT.md} \S2.1 (the five ingredients and the $522=2(K_{\max}+1)$ width identity),
\S2.2 (all intervals with widths and tails); \code{R03-CERT-THETA.md} \S2.1 (the two bugs, both found
by the $\theta=\pi$ regression, which reproduces \code{q11} to $4\times10^{-41}$), \S2.2 (the
interval table); \code{R10-AUDITCERT.md} \S2 (an independently written interval class and an
independent quadrature with a true remainder bound, not a heuristic one).}

\begin{claimbox}{cE}
\claimhead{Evidence}{(one number in this chapter had to be withdrawn, and the audit found it)}
\emph{The first certificate at $\theta\neq\pi$ reported $\mathrm{gap}(0)=0.729924127576750\ldots$
with an enclosure of width $1.4\times10^{-10}$; an independent implementation gives
$0.729924128620549$, higher by $1.044\times10^{-9}$, and that enclosure does \emph{not} contain the
true value. The four molecular integrals $Z,N_{A},N_{B},N_{S}$ and their tails \emph{are} valid
enclosures --- the audit verifies that its own $Z\,e^{-\beta c}$ lies inside them, at a distance
equal to their own tails ($7.4\times10^{-11}$ and $1.8\times10^{-21}$) --- so the defect is width loss
in the quotient: the intervals for $E[g]$ and $E[\Phi'']$ were printed with collapsed endpoints and
the tails did not propagate through the division. The same audit finds the $\theta=\pi/2$ enclosure
exceeded by $50\times$ for the same reason. The sign verdict never moved (the true value is
\emph{larger}), and a second route had already flagged the same $7.8\times10^{-10}$ discrepancy from a
different pair of coordinates.} \LE{}
\end{claimbox}

\src{\code{R10-AUDITCERT.md} \S0, \S3 (the diagnosis, the $7.8\times10^{-10}$ in $E[g]$ and the
$6.9\times10^{-11}$ in $E[\Phi'']$ whose weighted sum is the $1.04\times10^{-9}$ of the gap) and \S4
(that $D\equiv0$ at $\theta=0$ is an identity, so that clause is unaffected);
\code{R06-GAPMONO.md} \S4 (the independent flag, from $u$- and $x$-coordinate quadratures).
\textbf{What stands}: \eqref{eq:h1s-points}, with $\theta=0$ at $\pm10^{-55}$ rather than
$\pm10^{-10}$, and with the $\theta=0$ entry \emph{corrected} rather than inherited.}

\subsubsection*{From three points to the whole interval}

Three points do not make an interval. The passage uses one structural fact and one measurement, and
they are of different kinds, so the grading is split accordingly.

\emph{The structural fact} \LT. Put $c=\cos(\theta/2)$ and $x=\cos u$; then
$\sin^{2}u\,u^{-1}\mathrm{d} u=w(x)\mathrm{d} x$ with $w(x)=\sqrt{1-x^{2}}/\arccos x$, so
\begin{equation}
\nu_{c}\ \propto\ e^{\beta cx}\,w(x)\,\mathrm{d} x,\qquad x\in(-1,1),
\label{eq:h1s-expfam}
\end{equation}
a one-parameter exponential family with sufficient statistic $x$. Differentiating $\log Z(c)$ gives
the covariance formula $\partial_{c}E_{\nu}[f]=\beta\operatorname{Cov}_{\nu}(f,x)$ for every $f$, and
the derivative of the gap splits into five terms:
\begin{equation}
\partial_{c}\,\mathrm{gap}=\underbrace{\tfrac{\beta}{2}E[g]}_{T_{1}>0}
+\underbrace{\tfrac{\beta^{2}c}{2}\operatorname{Cov}(g,x)}_{T_{2}>0}
+\underbrace{\beta\operatorname{Cov}(\Phi'',x)}_{T_{3}<0}
+\underbrace{2\beta^{2}c\operatorname{Var}(g)}_{T_{4}>0}
+\underbrace{T_{5}}_{\text{sign open}} .
\label{eq:h1s-mono}
\end{equation}
$T_{1}>0$ is trivial; $T_{2}>0$ is the rearrangement inequality, since $g$ and $x$ are both strictly
decreasing in $u$; $T_{4}>0$ is the derivative of the explicit factor $-(1-c^{2})$; $T_{3}<0$ is the
\emph{only} negative term and it is negative for a reason that is itself provable, namely that
$\Phi''$ is strictly increasing on $(0,\pi)$ --- equivalently $F(u):=u^{2}+t^{2}+t-2>0$ with
$t=u\cot u$, and $F\sim2u^{4}/45$ near $0$. The fifth term collects the motion of the measure itself,
\begin{equation}
T_{5}=-\beta^{3}(1-c^{2})\bigl[\operatorname{Cov}(g^{2},x)-2E[g]\operatorname{Cov}(g,x)\bigr],
\label{eq:h1s-T5}
\end{equation}
and \emph{its sign is open}: it is the exact missing step. Measured, the positive terms dominate it
comfortably --- $\min(T_{1}-|T_{3}|)=0.3952$ and $\min(T_{2}+T_{4}+T_{5})=+0.0599$ over $51$ points ---
but a margin is not a proof.

\emph{The measurement} \LC. On $501$ interior points of $c\in(0,1)$ ($401$ uniform plus $101$
Chebyshev nodes at both ends, de-duplicated),
\begin{equation}
\partial_{c}\,\mathrm{gap}\in(0.44361829,\ 0.63898109),\qquad
\text{sign flips }0,\qquad\text{interior extremum }0 ,
\label{eq:h1s-501}
\end{equation}
and the analytic derivative agrees with a central difference to $2.9\times10^{-11}$ relative, so the
grid is measuring the identity \eqref{eq:h1s-mono} and not a discretization of it. Since
$\mathrm{d}\,\mathrm{gap}/\mathrm{d}\theta=-(s/2)\,\partial_{c}\mathrm{gap}\le0$ with equality only at the two
endpoints, the gap is strictly decreasing in $\theta$ and its unique worst point is $\theta=\pi$:
\begin{equation}
\min_{\theta\in[0,\pi]}\mathrm{gap}(\theta)=\mathrm{gap}(\pi)=0.1359947110\ldots>0 .
\label{eq:h1s-min}
\end{equation}

\begin{claimbox}{cC}
\claimhead{Conditional theorem}{(the whole interval, with the monotonicity as hypothesis)}
\emph{Assume \eqref{eq:h1s-expfam}--\eqref{eq:h1s-mono} (proved algebra) and the numerical sign
statement \eqref{eq:h1s-501} (not proved). Then $\mathrm{gap}(\theta)>0$ for every
$\theta\in[0,\pi]$ at $\beta=2.5$, with the quantitative lower bound \eqref{eq:h1s-min}. The three
certified points of \eqref{eq:h1s-points} are then \emph{redundant for the inclusion} and serve as
independent witnesses: monotonicity plus the single point $\theta=\pi$ already implies the interval,
and any reversal among the three would have falsified the monotonicity at once.} \LC{}
\end{claimbox}

\src{\code{R06-GAPMONO.md} \S1 (the exponential-family rewriting and the covariance formula), \S2 (the
five-term split and the numerical table of $T_{1},\dots,T_{5}$ at eleven values of $c$), \S3.1 (the
lemma $F>0$, i.e.\ $\Phi''$ strictly increasing, verified on $2000$ points with
$\min F=2.7\times10^{-13}$ consistent with $2u^{4}/45$), \S3.2 (\eqref{eq:h1s-501}, the grid, the
no-extremum result, and the finite-difference agreement). \textbf{The grade is deliberately
split}: the identities and $T_{1},T_{2},T_{4}>0$, $T_{3}<0$ are \LT; the sign of $T_{5}$ and hence the
interval statement is \LC, and this chapter writes it ``\LT{} modulo a \LC{} monotonicity'' nowhere
else without saying so.}

\emph{The cheap way to remove the hypothesis} \LC{}. A grid of certified points plus one Lipschitz
bound turns \eqref{eq:h1s-min} into an unconditional statement, and the budget is already computable
from \eqref{eq:h1s-501}: with $L:=\sup|\partial_{c}\mathrm{gap}|\le0.6390+|\mathrm{gap}''|\,\mathrm{d} c/2$
and $|\mathrm{gap}''|\le0.597$, a grid of spacing $\mathrm{d} c$ covers the interval as soon as
$L\,\mathrm{d} c<0.13599$. For $K=5$ ($\mathrm{d} c=0.200$) one gets $0.1397$, too large; for $K=6$
($\mathrm{d} c=0.167$) one gets $0.1148$, which succeeds; $K=8$ and $K=10$ give $0.0845$ and $0.0669$. Since
$c\in\{0,1/5,2/5,3/5,4/5,1\}$ makes $\beta c$ rational and the $\mathbb{Q}[\sqrt2]$ series machine
accepts rational data directly, \textbf{four additional certified points and one strict bound on
$|\mathrm{gap}''|$ make the whole of $[0,\pi]$ an unconditional \LT{} without any monotonicity
hypothesis.} That is compute, not mathematics, and it is the reason this chapter states the
conditional form rather than waiting for the proof of $T_{5}>0$.

\src{\code{R06-GAPMONO.md} \S0 (the honest label
``\LT{} $\mid$ monotonicity \LC'' and the two routes), \S6 (the $K$-budget, the numbers
$0.1397/0.1148/0.0845/0.0669$, and the remark that rational $c$ is directly acceptable), \S7 (what
is \LT{}, what is \LC{}, and the exact missing step).}

\emph{Two more strict fragments of the $\theta$ axis, for completeness} \LT. At $\theta=0$ the
statement holds for \emph{every} $\beta>0$ and not merely on a window: $s=0$ makes $G\equiv0$
pointwise, hence $D\equiv0$ identically in $(\beta,\nu)$, while $S\ge E[\Phi'']\ge1/6>0$. And there is
a purely analytic window needing no machine: with $\Phi''\ge1/6$, $0\le g\le1/2$ and hence
$\operatorname{Var}(g)\le1/16$, one has $S-D\ge1/6-\beta^{2}s^{2}/16>0$ whenever
$s^{2}<8/(3\beta^{2})$, i.e.\ $\theta\le0.4527\pi$ at $\beta=2.5$ --- a cruder bound than
\eqref{eq:h1s-min} but proved in three lines, with the endpoint $\theta=\pi$ supplied separately by
the exact structure $c=0$.

\src{\code{R03-CERT-THETA.md} \S3.2 items 2--3 (the analytic window $\theta\le0.4527\pi$, the
three-line derivation, and the $\forall\beta$ statement at $\theta=0$).}

\subsubsection*{What the certificate is not}

Four limits, and each one is a number rather than a caveat.

\textbf{(1) It is a window in $\beta$, and the window closes at a known coupling} \LT. At $\theta=\pi$
the value $S=1/(4Z)$ is independent of $\beta$ while $D=\beta^{2}V_{0}$ with
$V_{0}=0.0110592937729736$, so $S-D>0\iff\beta<\beta^{*}$ with
\begin{equation}
\beta^{*}=\sqrt{S/V_{0}}=4.3066078525597588749\ldots,
\label{eq:h1s-betastar}
\end{equation}
i.e.\ the certificate at $\beta=2.5$ has a $41.9\%$ margin, and at $\beta=4.4$ the same slice gives
$S-D=-0.00899<0$. \textbf{The certificate and the exact counterexample of
Chapter~\ref{ch:scalar-h1} are not in conflict and not in agreement: they live on the same slice at
different couplings.}

\textbf{(2) It does not extend to $a\neq0$} \LM. At $a\neq0$ the ratio $H/G$ is not constant
(\S\ref{sec:h1s-lemmas}, lemma 8), and what survives is an exact splitting with an explicit remainder,
not a criterion.

\textbf{(3) It does not extend to a larger volume} \LN. At $L=2$ there is no finite-dimensional
reduction at all (\S\ref{sec:h1s-limit}).

\textbf{(4) It is not a statement about Yang--Mills} \LN. Nothing in this subsection contains $a\to0$,
a continuum limit, or a coarse mode of a four-dimensional theory; $G''(0)$ here is a lattice-unit
curvature of a $0$-dimensional model at one coupling.

\src{\code{Q11-CERT.md} \S5 (the $\beta$ table $+0.2024/+0.1941/+0.1360/+0.0282/+0.0100/+0.00063/
-0.00899/-0.0714/-0.1930$ at $\beta=0.5,\dots,6$ and the $41.9\%$ margin), \S6 (that $L=2$ has no
exact integral representation and therefore no interval certificate);
\code{Q04-WEYL.md} \S4.3 (the same counterexample read as $E[H]\equiv0$ versus $\beta^{2}V_{0}>0$);
\code{ch03-existence.tex} \eqref{eq:dim} for why a lattice-unit curvature may not be read as a mass.}

\subsection{Eight strict lemmas, and what each one is allowed to say}
\label{sec:h1s-lemmas}

The certificate of \S\ref{sec:h1s-cert} is not an isolated computation: it stands on eight statements
that are proved rather than measured, and the value of the ledger is that each one is quoted with the
thing it does \emph{not} do. They are collected in Table~\ref{tab:h1s-lemmas} and then, because four of
them are routinely misquoted in the corpus, discussed one by one.

\footnotesize
\begin{longtable}{@{}>{\raggedright\arraybackslash}p{0.5cm}
  >{\raggedright\arraybackslash}p{6.4cm}
  >{\raggedright\arraybackslash}p{3.2cm}
  >{\raggedright\arraybackslash}p{4.0cm}@{}}
\caption{The eight strict lemmas of the day. ``Scope'' names the hypothesis that must travel with the
statement; a lemma quoted without its scope is false somewhere in this campaign's history.}
\label{tab:h1s-lemmas}\\
\toprule
\# & Statement & Source & Scope\\
\midrule
\endfirsthead
\toprule
\# & Statement & Source & Scope\\
\midrule
\endhead
\bottomrule
\endlastfoot
1 & $\Phi''(X;v)\ge\tfrac{N}{12}\langle v,v\rangle$ for every $X$ in the flat chart and every direction
$v$; sharp, attained only at $X=0$ & \code{P16-PHIQUANT.md} \S3, \S3.1 & all $N$, all measures, all
$\beta$; $\mathrm{SU}(N)$ with $\langle v,v\rangle=2\operatorname{tr}v^{2}$\\
2 & $E[H_{vv}\mid0]\ge0$, indeed $H_{vv}=\beta(c/2)g(u)\ge0$ pointwise & \code{P05-CINF.md} \S1.2 &
$\theta\le\pi$, this slice, $v=e_{3}$; false for $\theta>\pi$\\
3 & $E[H_{vv}]=E[G_{v}^{2}]$ and hence $E[H]-\operatorname{Var}(G)=E[G]^{2}\ge0$ & \code{Q15-SPECTRAL.md}
\S4 & measure $e^{-S_W}\mathrm{d}\mathrm{Haar}$, i.e.\ $\Phi$ absent; with $\Phi$ present a covariance term
remains\\
4 & $H=\tfrac12\cot(\theta/2)\,G$ pointwise & \code{Q04-WEYL.md} (2.2) & $a=0$ slice; $a\neq0$ has no
constant ratio\\
5 & $\operatorname{Var}(w)\le\tfrac12\langle w\rangle$, i.e.\ $2\operatorname{Var}\le\langle w\rangle$ &
\code{R02-CRIT.md} \S2; \code{R05-WCF.md} \S6.1; \code{R09-FLUCT.md} \S4 & unconditional in $(\beta,c)$;
three independent proofs\\
6 & $\langle w\rangle-2\operatorname{Var}(w)\ge2\langle w\rangle^{2}\ge0$ & \code{R05-WCF.md} \S6.1 &
unconditional; strictly stronger than 5\\
7 & $G''(0)=m_{G'}+E[\Phi'']-\operatorname{Var}(\Phi')-\operatorname{Cov}(G,\Phi')$ & \code{P21-REBUILD.md}
\S0 & algebra only; this is the \emph{skeleton} every other item is measured against\\
8 & $H=K_{0}G+\beta\Delta$ with $K_{0}=\tfrac12\cot(\theta/2)$ and $\Delta=[t\partial_{a}^{2}-c^{2}/(4t)](aw)$
& \code{R07-AFFINE.md} (4.1) & exact for all $a$; $\Delta$ is odd in $a$, first order $O(a)$\\
\end{longtable}
\normalsize

\paragraph{Lemma 1 (curvature floor).} The proof is two lines of calculus plus one convexity theorem.
With $\Phi=\sum_{i<j}g(\lambda_{i}-\lambda_{j})$, $g(d)=-2\log[\sin(d/2)/(d/2)]$, one has
$g''(d)=\tfrac12\csc^{2}(d/2)-2/d^{2}\ge1/6$; subtracting the quadratic part,
$h(X):=\Phi(X)-\tfrac{N}{12}\operatorname{tr}X^{2}=\sum_{i<j}\varphi(\lambda_{i}-\lambda_{j})$ with
$\varphi(d)=g(d)-d^{2}/12$, and $\sum_{i<j}(\lambda_{i}-\lambda_{j})^{2}=N\operatorname{tr}X^{2}$
makes the quadratic part of $h$ vanish; $\varphi''=g''-1/6\ge0$ makes $\varphi$ convex, hence $h$ is a
symmetric convex spectral function, hence convex in $X$ (Davis 1957), hence
$\operatorname{Hess}\Phi\succeq(N/6)\operatorname{tr}(\cdot)^{2}$, i.e.\
$\Phi''(X;v)\ge\tfrac{N}{12}\langle v,v\rangle$ after the normalization
$\langle v,v\rangle=2\operatorname{tr}v^{2}$. Because $\Phi$ is additive over links, the same line gives
$\Phi''\ge(N/12)\|v\|^{2}$ for \emph{every} direction, not merely for the rank-one ones --- this is the
step that earlier versions of the ledger claimed and did not have. Two honest qualifications: the
imported convexity theorem is checked at bibliographic level only, with an elementary hand proof for
$N=2$; and the floor is a floor, not a witness --- at $a=0.75$, $\beta=2.5$, $\theta=\pi$ one has
$\Phi''-\operatorname{Var}=+0.176>0$ while $G''(0)=-0.0059<0$, and the deficit is exactly
$E[H_{vv}]=-0.1816$. The floor cannot cancel the Wilson-side gap: $+0.167$ of floor plus at most
$+0.08$ of growth against $-0.947$ of $-\beta E[C_{\mathrm{geo}}]$. \LT{}~/~\LM

\src{\code{research/v3/species/P16-PHIQUANT.md} \S1 (the derivatives and the constant
$\min g''=0.166666675$), \S3 (the theorem, its two hypotheses, and the per-link additivity), \S3.1
(the numerical adjudication: $\inf\min\mathrm{eig}(\operatorname{Hess}\Phi)$ equals $N/12$ for
$N=2,\dots,6$ with ratio $1.000000$, attained at $X=0$), \S6 (why the floor does not cancel
$-\beta E[C_{\mathrm{geo}}]$, and the $a_{0}=0.75$ counterexample that keeps this lemma from being a
proof of $H1$). The Davis reference is C.~Davis, Arch. Math. \textbf{8} (1957) 276--278,
DOI 10.1007/BF01898787, cited at bibliographic level: the text was not read.}

\paragraph{Lemma 2 (sign of the mean curvature).} On this slice $H_{vv}/\beta=Q=\cos(\theta/2)g(u)/2$
with $g>0$ on $(0,\pi)$ and $g(0)=1/2$, so $Q\ge0\iff\cos(\theta/2)\ge0$. Hence the mean-side leg is
\emph{free} for $\theta\le\pi$, and $E[H_{vv}]=\beta(c/2)G(\beta c)$ with
$G(B)=E_{\nu_{B}}[g]\in(0,1/2)$ strictly. Beyond $\theta=\pi$ the statement reverses and the growth is
not linear: a boundary-layer expansion gives $E[H_{vv}]\simeq-0.12699\sqrt{|c|\beta}$, with measured
slopes $0.460/0.468/0.472$ at $\theta/\pi=1.25/1.5/1.75$. The lemma is therefore quoted with
$\theta\le\pi$ attached, and the corpus statement that $E[H_{vv}]$ is indefinite in general
($\min=-1.893$ on four-link random directions) is a statement about a different object, not a
contradiction. \LT{}~/~\LC

\src{\code{P05-CINF.md} \S1.2 (the pointwise computation and $G<1/2$), \S4 (the clause table), \S6.1
and \S6.2 (the $\theta>\pi$ branch, the $\sqrt\beta$ law, and the $0.6\%$ numerical check at
$\beta=10^{4}$); the four-link indefiniteness is \code{O06-CURV.md} and \code{LOOP-RUNBOOK.md}:2426 as
quoted by \code{P05-CINF.md} \S1.1.}

\paragraph{Lemma 3 (Stein, and why the spectral route is empty).} Take a left-invariant direction
field $v$; left invariance of Haar makes $\operatorname{div}_{\mathrm{Haar}}v=0$, so
$\int vf\,\mathrm{d}\mathrm{Haar}=0$ for all smooth $f$. Applying this to $f\,e^{-S_W}$ gives
$E_{\mu}[vf]=E_{\mu}[fG_{v}]$ with $\mu\propto e^{-S_W}\mathrm{d}\mathrm{Haar}$, and taking $f=G_{v}$ gives
$E[H_{vv}]=E[G_{v}^{2}]$; hence $E[H]-\operatorname{Var}(G)=E[G]^{2}\ge0$ with equality iff $E[G]=0$.
On a single plaquette the equality is exact for every $\beta>0$ (the $n\mapsto-n$ symmetry kills
$E[G]$), and the measured value at $\beta=2.5$ is $E[H]=\operatorname{Var}(G)=1.26799$ with a
difference of $6\times10^{-7}$. The corollary is a no-go and belongs here rather than in
\S\ref{sec:h1s-roads}: the spectrum of the Hessian enters the difference only through
$E[\lambda_{k}]=E[g_{k}^{2}]$ and cancels term by term, so no spectral condition is equivalent to it;
and the one named spectral equivalence that does exist, $\Lambda\preceq I$ (Brascamp--Lieb /
$CD(0,\infty)$), is \emph{false} here --- $\lambda_{k}=\beta\cos\varphi$ threefold degenerate,
$\lambda_{\max}=\beta=2.5>1$, $P(\lambda<0)=0.105$ --- while the target inequality holds and is an
equality. \LT{}/~\LN

\src{\code{Q15-SPECTRAL.md} \S1--\S2 (the identities), \S3 (the spectral table and
$P(\lambda\gtrless)$), \S4 (the proof by the divergence identity, and the statement that a spectral
proof is unavailable).}

\paragraph{Lemma 4 (a pointwise ratio, which is why the criterion is one inequality).} On the $a=0$
slice $G=\beta s\,w(\rho)$ and $H=\tfrac{\beta c}{2}w(\rho)$ with $w(\rho)=\sin(\rho/2)/\rho$, so
$H=\tfrac12\cot(\theta/2)G$ \emph{pointwise}, with a measured deviation of $8.9\times10^{-16}$ and
finite-difference agreement to $1.2\times10^{-8}$. This is an equality, not an inequality, and its
consequence is that \emph{every} attempt at a pointwise comparison $H(\rho)\ge G(\rho)^{2}/\cdots$ is
empty: the comparison collapses to one number,
\begin{equation}
\operatorname{Var}(G)\le E[H]\iff 2\beta s^{2}\operatorname{Var}_{\mu}(w)\le c\,\langle w\rangle ,
\label{eq:h1s-crit}
\end{equation}
which is the criterion of \S\ref{sec:h1s-window}. Off $a=0$ no such collapse exists because $H$ and
$G$ stop being proportional --- and \S\ref{sec:h1s-lemmas} lemma 8 is the exact measure of that
failure. \LT

\src{\code{Q04-WEYL.md} (2.1)--(2.2) and \S4.1 (``the pointwise content is an equality, not an
inequality''); \code{R09-FLUCT.md} \S1--\S2 (the same pair read as $\partial_{a}h$ and
$\partial_{a}^{2}h$, with $S_W=\beta h$).}

\paragraph{Lemmas 5 and 6 (two moment inequalities, three independent proofs).} Both are statements
about the single function $w$, not about the model, and both are unconditional.

\emph{Path (a)}, from the kernel: $A-2B=\tfrac12\int e\,\sin^{3}s\,(s-\sin s)/s^{3}\,\mathrm{d} s$ where
$e=e^{-\beta(1-c\cos s)}$. The kernel is positive because $s-\sin s>0$, and strictly decreasing;
pairing the humps $s_{\mp}=2\pi k+\pi\mp u$ around each zero of $\sin$ gives exactly equal exponents
and exactly equal $|\sin^{3}|$, so $g(s_{-})>g(s_{+})$ and each pair contributes positively, with the
final partial hump also positive because $\Lambda=11\pi$ ends on the positive side. Hence
$2\langle w^{2}\rangle\le\langle w\rangle$ and therefore $\operatorname{Var}(w)\le\tfrac12\langle w\rangle$
for all $\beta\ge0$, all $c\in[0,1]$, and on the multi-period window $\Lambda=22\pi$ where $w$
changes sign. Sharpness: $2\langle w^{2}\rangle/\langle w\rangle\to1$ as $\beta\to\infty$
(mass concentrates at $s\to0$, $w\to1/2$), measured maximum $0.996792$, so the constant cannot be
improved in that form. A negative result belongs to this path and prevents a wrong proof:
$E|w|\le\langle w\rangle$ is \emph{false} ($E|w|/\langle w\rangle=1.197947$ at $\Lambda=22\pi$,
$\beta=0$), because $\int|w|$ adds the paired contributions that $\int w$ subtracts; any route through
$|w|\le1/2\Rightarrow\operatorname{Var}\le\tfrac14\langle w\rangle$ is therefore closed.

\emph{Path (b)}, five lines and measure-free: $w=\tfrac12-D$ with $0\le D\le\tfrac12$, so
$D^{2}\le D/2$ and
\begin{equation}
\langle w\rangle-2\operatorname{Var}(w)
=\bigl(\tfrac12-\langle D\rangle\bigr)-2\bigl(\langle D^{2}\rangle-\langle D\rangle^{2}\bigr)
\ge\tfrac12-2\langle D\rangle+2\langle D\rangle^{2}
=2\langle w\rangle^{2}\ge0 .
\label{eq:h1s-r05}
\end{equation}
This is lemma 6, and it is strictly stronger than lemma 5, equivalently
$\operatorname{Var}(w)\le\tfrac12\langle w\rangle(1-2\langle w\rangle)$.

\emph{Path (c)}, from the same kernel in the $u$-coordinate:
$E[g]-2E[g^{2}]=(1/2Z)\int e\,\sin^{3}u\,(u-\sin u)/u^{3}\,\mathrm{d} u\ge0$ pointwise, verified digit by
digit against $\mathrm{Var}$ on fifteen $(\beta,c)$ pairs (e.g.\ $\beta=2.5$, $\theta=\pi$:
$7.984483\times10^{-3}$ on both sides), with a pointwise kernel minimum $\ge1.9\times10^{-22}>0$.

All three paths give $\operatorname{Var}(w)\le\tfrac12\langle w\rangle$; they were derived from
different objects ($A-2B$ versus $D^{2}\le D/2$ versus the direct kernel identity) and none of them
uses the sign of $\beta c$ or the value of $\theta$. \LT

\src{\code{R02-CRIT.md} \S2 (path (a), including the pairing argument, the sharpness, and the
negative result $E|w|\le\langle w\rangle$ false); \code{R05-WCF.md} \S6.1 (path (b) verbatim, and the
comment that lifting the numerical factor $16.16$ requires the missing monotonicity);
\code{R09-FLUCT.md} \S4 (path (c), the identity, and the $15$-point digit check);
\code{LOOP-RUNBOOK.md} \S11.170 (\textbf{the explicit statement that these are three independent
paths to one conclusion}).}

\paragraph{Lemma 7 (the skeleton).} With $F=\partial_{v}V=G+\Phi'$ and $\kappa=\partial_{v}^{2}V=H+\Phi''$,
\begin{equation}
\operatorname{Var}(F)-E[\kappa]=-G''(0)=-m_{G'}-E[\Phi'']+\operatorname{Var}(\Phi')
+\operatorname{Cov}(G,\Phi'),
\label{eq:h1s-skeleton}
\end{equation}
which is algebra: it combines $\Delta_{\mathrm{dyn}}=m_{G'}+\operatorname{Cov}(G,\Phi')$,
$G''(0)=\Delta_{\mathrm{dyn}}+W_{\mathrm{geo}}+X$, $W_{\mathrm{geo}}=E[\Phi'']-\operatorname{Var}(\Phi')$
and $X=-2\operatorname{Cov}$. Its numerical content is a check, and it checks twice: on the six-point
$a$-grid the residual is $\le7\times10^{-6}$, and on $L=2$ at $\beta=2.5$ with the archived values
$0.562+1.09+0.043=1.695$ against the measured $1.696\pm0.293$, a $0.06\%$ deviation. The corpus
contains an acid test that this identity also passes and that reorganizes the tool set: of the four
candidate quantities $\{\operatorname{Cov}(G,\Phi'),\operatorname{Var}(\Phi'),m_{G'},
E[\Phi'']-N/12\}$, the three that involve $\Phi$ alone \emph{fail} the free-field test ---
$\operatorname{Var}(\Phi')$ tends to $2.341\times10^{-3}$, which is $65\%$ of its interacting value
rather than zero, and $E[\Phi'']-N/12$ is \emph{identical} at $\beta=2.5$ and $\beta=10^{-6}$ because
at $\theta=\pi$ the measure is Haar for every $\beta$. What survives is $\{m_{G'},-\operatorname{Cov}\}$,
both containing $G=\partial_{v}S_{W}$, the only object carrying a power of $\beta$. \LT{}/~\LE

\src{\code{P21-REBUILD.md} \S0 (the identity, its numerical check, and the $L=2$ closure to
$0.06\%$), \S1--\S2 (the four candidates, the acid test, the free-field values, and the reorganized
pair), \S3 (the scope list); the split $G''=\Delta_{\mathrm{dyn}}+W_{\mathrm{geo}}+X$ is
\code{ch05-theorems.tex} \eqref{eq:split}.}

\paragraph{Lemma 8 (an exact splitting instead of an approximate affinity).} The corpus had been
using $H=K(a)G+C(a)$ with $R^{2}>0.9992$. That is not an identity, and the proof that it is not is
worth recording because the check is a Wronskian rather than a fit: $W(G,H)=GH'-G'H\not\equiv0$ and
$W(G',H')=G'H''-G''H'\not\equiv0$, both of first order $O(a)$ (numerically $-7.21\times10^{-4}$ and
$-8.46\times10^{-6}$ at $a=0.3$, $\rho=2.5$); and $H-KG$ has no constant term in the basis
$\{\sin(\rho/2),\cos(\rho/2),1\}$ since $W\{\sin,\cos,1\}=-1/8\neq0$, so a strict affinity would force
$C=0$, i.e.\ a strict proportionality --- which the ratio computation kills as well
($Q_{H}/Q_{G}=K_{0}+3/a-3a/\rho^{2}$ varies with $\rho$, and the two channels differ by a factor
$30$: $K_{c}=10.356$ against $K_{s}=0.308$ at $a=0.3$). What \emph{is} exact is
\begin{equation}
H=K_{0}G+\beta\Delta,\qquad K_{0}=\tfrac12\cot(\theta/2),\qquad
\Delta=\Bigl[t\,\partial_{a}^{2}-\frac{c^{2}}{4t}\Bigr]\bigl(a\,w(\rho)\bigr),
\label{eq:h1s-affine}
\end{equation}
with $\Delta$ odd in $a$ and of first order, and hence $K_{\mathrm{LS}}=K_{0}+\beta\operatorname{Cov}(G,\Delta)/
\operatorname{Var}(G)$, $C_{\mathrm{LS}}=\beta[E\Delta-(\operatorname{Cov}/\operatorname{Var})E[G]]$ and
$1-R^{2}=\beta^{2}\operatorname{Var}(\Delta)(1-\mathrm{corr}^{2})/\operatorname{Var}(H)$: the fit
residual is not noise but a computed variance, and it is of order $O(a)$ with
$\mathrm{relres}/a\to1.43\times10^{-3}$, invariant under raising the quadrature from $1500$ to $6000$
points. \LT{}/~\LE

\src{\code{R07-AFFINE.md} \S2 (the three exact identities and the mechanism: one operator $L=t\partial_{a}+c/2$
acting on two adjacent $a$-derivatives), \S3 (the two Wronskians and the $C=0$ argument), \S4--\S5
(the splitting, the closed forms for $K$ and $C$ reproducing the fit to $10^{-16}$, and the residual
table).}

\subsection{One window, at two grades}
\label{sec:h1s-window}

Combining lemmas 4 and 5 turns \eqref{eq:h1s-crit} into a region of $(\beta,\theta)$ in which H1
holds with no measurement at all, and the region has a strictly provable boundary and a measured one.

\begin{claimbox}{cT}
\claimhead{Theorem}{(the analytic window)}
\emph{With $w(\rho)=\sin(\rho/2)/\rho$, $\mu\propto e^{-\beta(1-c\cos(\rho/2))}\sin^{2}(\rho/2)\rho^{-1}\mathrm{d}\rho$
on the $a=0$ slice, \eqref{eq:h1s-crit} holds whenever}
\begin{equation}
\beta\sin^{2}(\theta/2)\ \le\ \cos(\theta/2),
\qquad\text{i.e.}\qquad
\beta\le\frac{\cos(\theta/2)}{\sin^{2}(\theta/2)} .
\label{eq:h1s-wintight}
\end{equation}
\emph{Consequently $H1$ holds for every $\beta\ge\beta_{0}(\theta):=\sin^{2}(\theta/2)/\cos(\theta/2)$;
at $\beta=2.5$ this covers $\theta\le0.388\pi$ (from $c\ge0.8198$) and at $\beta=6$ it covers
$\theta\le0.256\pi$ (from $c\ge0.9201$).} \LT{}
\end{claimbox}

\src{\code{Q04-WEYL.md} \S4.1 (4.1)--(4.2): the implication from $0\le w\le1/2$ through
$\operatorname{Var}(w)\le\tfrac12\langle w\rangle$ to \eqref{eq:h1s-wintight}, with the numerical
coverage $\theta/\pi\le0.7281/0.5758/0.3881$ at $\beta=0.5/1/2.5$; \code{R02-CRIT.md} \S3 (the same
implication, now with the kernel proof of the moment inequality, so that the step is a theorem and not
an observation); \code{R09-FLUCT.md} \S4 (the third derivation, ending at the same threshold and
noting $\theta/\pi\le0.3882$ at $\beta=2.5$).}

The measured window is wider, and the gap between the two grades is exactly the factor that lemma 6
buys. From \eqref{eq:h1s-crit} and lemma 6 one may replace $\tfrac12\langle w\rangle$ by
$\tfrac12\langle w\rangle(1-2\langle w\rangle)$, i.e.\ replace the constant $2$ in
$R:=\langle w\rangle/\operatorname{Var}(w)\ge2$ by the measured $R_{\min}$; the smallest value over
$\gamma=\beta c\in[0,\infty)$ is
\begin{equation}
R_{\min}=\frac{\langle w\rangle}{\operatorname{Var}(w)}\Big|_{\gamma=0}
=\frac{2N_{1}Z}{N_{2}Z-N_{1}^{2}}=32.3177375272 ,
\qquad Z=\tfrac12\mathrm{Cin}(2\pi),\quad
N_{1}=\tfrac34[\mathrm{Cin}(3\pi)-\mathrm{Cin}(\pi)],\quad
N_{2}=\mathrm{Cin}(4\pi)-\mathrm{Cin}(2\pi),
\label{eq:h1s-Rmin}
\end{equation}
so the criterion becomes
\begin{equation}
\beta\sin^{2}(\theta/2)\ \le\ \tfrac12\cos(\theta/2)\,R
\ \Longrightarrow\
\beta\le\beta_{c}(\theta)=16.15887\,\frac{\cos(\theta/2)}{\sin^{2}(\theta/2)}
\qquad\text{(measured tier)},
\label{eq:h1s-winc}
\end{equation}
which at $\beta=2.5$ gives $\theta\le0.904\pi$ and at $\beta=6$ gives $\theta\le0.786\pi$ --- a
$2.3$-fold widening in $\theta$ over \eqref{eq:h1s-wintight} at $\beta=2.5$ and a $3.1$-fold widening
at $\beta=6$. The corresponding
``$R\ge2$ but nothing better'' version, using only the proved lower bound $R_{\min}\ge4.017$ on the
multi-period window, gives $\theta\le0.527\pi$ and $\theta\le0.356\pi$: a $36\%$--$39\%$ widening that
is \emph{also} at the measured tier, because $R_{\min}=4.017$ there is a numerical minimum and not a
theorem.

\begin{claimbox}{cE}
\claimhead{Evidence}{(where the two grades meet, and why the coincidence is not a derivation)}
\emph{The measured boundary at $\beta=6$, $\theta\le0.786\pi$, sits within $0.9\%$ of a boundary
obtained from a different functional by a different criterion:
$\beta_{0}(\theta)=\tan^{2}(\theta/2)/(36c_{\infty})$ with $c_{\infty}=\cos(\theta/2)/4$ gives
$\theta_{\max}=0.7792\pi$ at $\beta=2.5$ and $0.8335\pi$ at $\beta=6$. The two routes do not compute
the same thing --- \eqref{eq:h1s-winc} tests $\Delta_{\mathrm{dyn}}\ge0$, i.e.\ $\operatorname{Var}\le E[H]$
alone, whereas the $\beta_{0}$ route tests $G''(0)\ge0$ with $W_{\mathrm{full}}$ included --- so
agreement is not required and the $0.9\%$ should be read as evidence that both are describing the
geometry of the same slice, not as a proof of either.} \LE{}
\end{claimbox}

\src{\code{R05-WCF.md} \S2--\S5 (\eqref{eq:h1s-Rmin}: the closed forms at $\beta=0$, the $1/\gamma$
expansion for $\beta>0$, the $41$-point table, and the minimum at $\gamma=0$, which is a boundary
value and hence exact in $\gamma$), \S6.2--\S6.3 (\eqref{eq:h1s-winc}, the factor $16.158868764$,
and the sentence that the missing step is the monotonicity of $\operatorname{Var}(w)$ in $\gamma$);
\code{LOOP-RUNBOOK.md} \S11.169 (the two windows side by side, $0.904\pi$ and $0.786\pi$, and the
$0.786\pi$ versus $0.7791\pi$ coincidence); \code{P20-THETA.md} \S4.2 (the $\beta_{0}$ route, its
anchor values $2.50$ at $0.7792\pi$ and $6.00$ at $0.8335\pi$, and $\S2$ for why $\beta_{0}$ is a
ratio of a $\beta$-free $V_{\infty}$ to a $\beta$-free $c_{\infty}$).}

Two boundaries must be quoted with the window, not after it.

\textbf{(1) No sharpening of these two legs makes the window $\theta$-uniform} \LN. The criterion is a
product, $\beta\,w(\theta)\ge\tfrac19$ with $w(\theta)=\cos^{3}(\theta/2)/\sin^{2}(\theta/2)$, and
$w(\pi)=0$: hence $\sup_{\theta}\beta_{0}(\theta)=+\infty$, and no injective reparametrization of
$\theta$ can change that, because the product form is preserved. Quantitatively
\begin{equation}
\beta_{0}(\theta)\ \asymp\ \frac{8}{9(\pi-\theta)^{3}}\qquad(\theta\to\pi),
\label{eq:h1s-div}
\end{equation}
which reproduces the tabulated $2.87\times10^{4}$ at $\theta=0.99\pi$ to three digits. The escape is
not a better constant on the same legs but a different decomposition --- and the $\theta=\pi$ slice is
where the decomposition fails, not where it is loose.

\textbf{(2) At $\theta=\pi$ the premises of both legs collapse} \LM. There $c=0$, so $B=\beta c\equiv0$,
the family $\nu_{B}$ degenerates to the single Haar law, $\operatorname{Var}=\beta^{2}V_{0}$ is
unbounded in $\beta$, $V_{\infty}=+\infty$ (the variance leg is vacuous, not false), $c_{\infty}=0$ and
$E[H_{vv}]\equiv0$ (the mean-side leg is vacuous too). This is exactly the slice on which
\eqref{eq:h1s-betastar} is exact, so the certificate and the counterexample are the two faces of one
degeneration. There is, however, one honest widening on the other side: a small-$\beta$ leg,
$\operatorname{Var}\le\beta^{2}V_{0}$ with $V_{0}=W(0)$, gives $H1$ for \emph{every} $\theta$ provided
$\beta^{2}V_{0}+\delta\le\inf W_{\mathrm{full}}$, i.e.\ $\beta\lesssim3.9$ --- at the measured tier,
because $\inf W_{\mathrm{full}}$ is a table infimum. What remains uncovered is therefore the bounded
rectangle $\{\theta>0.8335\pi\}\times\{\beta\in(3.9,6]\}$, and that rectangle contains the known false
point $(\theta=\pi,\beta>4.30661)$: it is not ``unproved'' but \emph{false}.

\src{\code{P20-THETA.md} \S4.1 (the product form, the impossibility, and the observation that only a
different decomposition can help), \S2.3 (why $B=\beta c$ is not a substitute variable: the threshold
goes to zero in $B$-language while the inequality becomes $\beta\le\beta^{*}$), \S3 (the premise
collapse table, each row \LT{} or \LC{} as marked), \S4.4--\S4.5 (the small-$\beta$ leg, $\beta_{\rm
unif}\approx3.9$, the residual rectangle, and that it contains a false point);
\code{O22-THRESHOLD.md} \S0 and \S5 as quoted there for the numerical anchors; \code{Q11-CERT.md} \S5
for \eqref{eq:h1s-betastar} on the same slice.}

\subsection{The lattice ledger at $L=2$ and $L=3$}
\label{sec:h1s-lattice}

Everything above is a $0$-dimensional statement. The measured side of the ledger is a Monte Carlo on
the four-dimensional model at $L=2$ and $L=3$, and its content is the sign of the same quantity,
$G''(0)=E[\kappa]-\operatorname{Var}(F)$, with the split observation $S_W$ versus $\Phi$ maintained per
sample so that $\Delta_{\mathrm{dyn}}+W_{\mathrm{geo}}+X-G''(0)$ can be checked on every configuration
(the sum check is $\le1.6\times10^{-6}$ at $\beta=6$ and $\le1.6\times10^{-5}$ across the $\beta$
scan).

\begin{center}
\footnotesize
\begin{tabular}{@{}lll>{\raggedright\arraybackslash}p{6.4cm}@{}}
\toprule
volume & $\beta$ & $G''(0)$ & statistics\\
\midrule
$L=2$ & $1$   & $+0.2835\pm0.0049$ (jk) & $27.0\sigma$ (jk)\\
$L=2$ & $2.5$ & $+0.4116\pm0.0421$ (jk) & $4.0\sigma$ (jk); note also $+0.3942\pm0.0258$ at $6.0\sigma$
from a $1.5\times10^{4}$-sweep run\\
$L=2$ & $4$   & $+0.5108\pm0.0939$ (jk) & $2.3\sigma$ (jk)\\
$L=2$ & $6$   & $+0.6018\pm0.0589$ (jk25) & $10.2\sigma$; five estimators span $4.1$--$10.2\sigma$\\
$L=3$ & $2.5$ & $+0.3396\pm0.0585$ (jk) & $5.80\sigma$ (jk) / $2.24\sigma$ ($\tau$)\\
\bottomrule
\end{tabular}
\end{center}

Three readings, in increasing order of what they cost.

\emph{The sign is uniform in $\beta$ on this window} \LE. At $L=2$ the four central values are
monotone increasing, $0.2835\to0.4116\to0.5108\to0.5913$, with no interior minimum and no crossing;
$16/16$ chains have positive central values, the weakest being $+0.1338$ ($\beta=4$) and $+0.2258$
($\beta=6$). The mechanism is a competition that does not flip: $D/S$ rises $0.505\to0.787\to0.848\to
0.889$, i.e.\ it approaches $1$ without crossing it, so $G''(0)=E[\kappa](1-D/S)$ stays positive while
its margin shrinks --- a fit gives $G''\propto\beta^{0.41}$. \textbf{The extrapolation beyond
$\beta=6$ is an extrapolation}: a quadratic fit puts the zero near $\beta\approx17.3$ and a cubic near
$14.9$, and those numbers are quoted here as extrapolations, not as results.

\emph{The statistics are enough at $\beta=6$ only after the estimator is fixed} \LE. The first
$L=2$ determination at $\beta=6$ ($4$ chains $\times\,5\times10^{3}$ sweeps) gave $+0.5913$ but only
$1.3\sigma$ under the $\tau$-corrected error, because $\tau_{K}$ grows to $100$--$240$ and $N_{\rm
eff}$ falls to $21$--$71$. Re-running with $8$ chains $\times\,2\times10^{4}$ sweeps (burn
$2\times10^{3}$, seeds $700$--$707$, $N_{\rm eff}=48$--$194$ per chain, $818$ in total) gives
$+0.6018$ with five estimators:
jk25 $\sigma=0.0589$ ($10.2\sigma$), eight-chain scatter $0.0651$ ($9.2\sigma$), long-block
batch means $0.0733$ ($8.2\sigma$), $\tau_{z}$-corrected $0.0823$ ($7.3\sigma$), and the old
$\tau_{F}$ formula $0.1450$ ($4.1\sigma$). $8/8$ chains are positive (sign test $p=3.9\times10^{-3}$),
the half-split difference is $-0.2\sigma$, $\chi^{2}/\mathrm{dof}=1.22$. The one estimator still below
$5\sigma$ is the old one, and the audit of it is a formula mismatch rather than a matter of taste: it
applies $\tau_{F}$ to the fluctuation of the sample variance $z=(F-\bar F)^{2}$, whose correlation time
is $\tau_{z}=8.7$--$14.1$, not $\tau_{F}=24$--$48$, and it treats a fourth moment as Gaussian
($\mathrm{Var}(z)=2\sigma^{4}$); with $\tau_{z}$ the same machinery returns $7.3\sigma$. \textbf{Even
under the old estimator the requirement is explicit}: $2.9\times10^{4}$ sweeps per chain, $1.45$ times
the present run.

\emph{There is no volume dependence} \LE. At $\beta=2.5$ the two volumes give $+0.3942$ ($L=2$) and
$+0.3396$ ($L=3$), a difference of $-0.055\pm0.064$, i.e.\ $0.85\sigma$; and the quantity that would
carry a volume effect, $\operatorname{Var}(F)$, agrees to $0.7\%$ ($1.535$ versus $1.5458$). The $L=3$
number also settles a bug: the earlier negative value $-0.177\pm0.062$ came from a sampler with no
fundamental-domain constraint, which let the chain wander onto the repeated flat chart where
$\Phi\sim2\log(r/2)$ supplies no restoring force; with the hard rejection $|x_{l}|<2\pi$ (rejection
rate $0.14\%$) $\operatorname{Var}(F)$ falls from $1.8917$ to $1.5458$ and the sign flips, exactly as
the domain analysis predicted.

\src{\code{P19-BETA-AT-L2.md} \S0 and \S2 (the four central values, $16/16$ chains, the monotone
trend, the $D/S$ mechanism, and the explicit label on the extrapolations);
\code{Q13-BETA6.md} \S0--\S3 (the $\beta=6$ value, the six estimators, $8/8$ chains, the sign test,
$\chi^{2}/\mathrm{dof}$, the $\tau_{z}$ analysis and the $2.9\times10^{4}$-sweep requirement);
\code{P01-L3.md} \S0 and \S1 (the $L=3$ value, $5.80\sigma$/ $2.24\sigma$, the fundamental-domain
bug and its mechanism); \code{O24-L2-PUSH.md} \S2 and \S5 (the $1.5\times10^{4}$-sweep
$+0.3942\pm0.0258$ at $6.0\sigma$, the two error estimators, and the limits of the run);
\code{LOOP-RUNBOOK.md} \S11.160--\S11.161 (the combined statement and the $0.85\sigma$/$0.7\%$
comparison). \textbf{Every entry is \LE: one volume window, one coarse mode (A), no $a\to0$, and a
finite-difference definition of $\kappa$ and $F$ with $h=10^{-3}$.}}

\subsection{Twelve roads that end, each with its number}
\label{sec:h1s-roads}

The campaign's negative asset is a list of routes that cannot work, each ended by a number rather than
by an impression. Twelve are in force after this day: ten that were already in force, and two that
this day's work added --- the curvature-fluctuation form, which turned out to be the criterion renamed
and dies at an exact point, and the affine reduction, whose residual is not noise.

\footnotesize
\begin{longtable}{@{}>{\raggedright\arraybackslash}p{0.4cm}
  >{\raggedright\arraybackslash}p{3.2cm}
  >{\raggedright\arraybackslash}p{8.2cm}
  >{\raggedright\arraybackslash}p{2.8cm}@{}}
\caption{The twelve roads, each with the number that ends it. All twelve are class-level
\LN: none is ``the hypothesis was not established'', and in every row but two the hypothesis is
\emph{false}.}\label{tab:h1s-roads}\\
\toprule
\# & Road & Why it dies, with the number & Source\\
\midrule
\endfirsthead
\toprule
\# & Road & Why it dies, with the number & Source\\
\midrule
\endhead
\bottomrule
\endlastfoot
1 & Brascamp--Lieb, all directions & the required $\lambda$ is $38$--$44$ times $N/12$, and
$\operatorname{Hess}(S_{W}+\Phi)$ is indefinite at $\beta=2.5$: the hypothesis fails, not the
constant & \code{Q02-BL-PHI.md}\\
2 & ODE / Riccati for $G''$ & $G''$ does not close: $D'\le S'$ is \emph{equivalent} to $G'\ge0$
($\max|\mathrm{gap}+G'|=2\times10^{-16}$), i.e.\ a tautology that cannot be upgraded; and $D/S$ crosses
$1$ at $\beta_{c}\approx8$--$11$, with $\beta=6$ already at $0.941$ & \code{Q06-RICCATI.md}\\
3 & Curvature-dimension, RC-BL & $\mathrm{ess\,inf}\,K_{+}=0$ (the bad set carries $K_{+}\equiv0$;
$K\ge N/3-|\lambda_{\min}|=-0.893$ for $\mathrm{SU}(3)$), and $\sum\theta_{j}$ cannot be made
consistent & \code{Q14-RCBL.md}\\
4 & Sum of squares & the direction-by-direction pairing needs
$\operatorname{Var}\le N/12-\langle|\lambda_{-}|\rangle=-0.0367<0$ with $\langle\lambda_{-}\rangle=-0.2033$:
false for every $\operatorname{Var}\ge0$, pure arithmetic & \code{Q12-SOS.md}\\
5 & Cauchy--Schwarz & loose by $2.6\times$ on the upper side and by $40$--$125\times$ on the lower side &
\code{Q09-COVLOW.md}\\
6 & Spectral conditions & $\lambda_{k}$ cancels term by term in $E[H]-\operatorname{Var}(G)$; the
named spectral criterion $\Lambda\preceq I$ holds on $0.00\%$ of configurations
($\lambda_{\max}=\beta=2.5$) while the target holds and is an equality & \code{Q15-SPECTRAL.md}\\
7 & Two-halves decomposition & the $\Phi$-only margin is \emph{exactly} $0$ unconditionally
($1/3$ against $1/3$, by parts at machine precision), and on the interacting side $\Phi$ turns negative
at $a>1.588$ & \code{Q01-PHI-ONLY.md}\\
8 & The $\beta_{0}$ window's small-$B$ bridge & the claimed theorem-grade bound
$\operatorname{Var}\le\beta^{2}s^{2}(1/2-1/\pi)^{2}/4$ is broken by a counterexample at its own
window: at $B=0.5$, $\Psi=0.082919>0.2971B^{2}=0.0743$ & \code{P20-THETA.md} \S6.5\\
9 & $\Phi$-only tools as the witness & they do not vanish on the free field:
$\operatorname{Var}(\Phi')=2.341\times10^{-3}$, which is $65\%$ of the interacting value, and
$E[\Phi'']-N/12$ is identical at $\beta=2.5$ and $10^{-6}$ & \code{P21-REBUILD.md}\\
10 & Statements uniform in $a$ & the negative region grows monotonically in $\beta$: width
$1.09\to2.26\to3.57\to4.36\to5.22\to5.66$ and depth $-0.016\to-22.7$ for $\beta=2\to40$ & \code{Q08-LARGEA.md}\\
11 & Curvature fluctuation $\beta\operatorname{Var}(\partial_{t}h)\le E[\partial_{t}^{2}h]$ &
it is \eqref{eq:h1s-crit} renamed ($S_W=\beta h$ makes both sides the same ratio) and fails at
$\theta=\pi$ for every $\beta>0$: LHS $=0.027648>0=$ RHS, because $H_{vv}\equiv0$ is a pointwise
identity & \code{R09-FLUCT.md}\\
12 & Affine reduction $H=K(a)G+C(a)$ & not an identity: $W(G,H)=O(a)\not\equiv0$,
$W(G',H')\not\equiv0$, the two channels differ $30$-fold ($K_{c}=10.356$ against $K_{s}=0.308$), and the
residual is $O(a)$ with $\mathrm{relres}/a\to1.43\times10^{-3}$, stable under a fourfold quadrature
increase & \code{R07-AFFINE.md}\\
\end{longtable}
\normalsize

\begin{claimbox}{cN}
\claimhead{No-go}{(the classical variance--curvature toolset, as a class)}
\emph{Rows 1--6 exhaust the classical families that produce ``variance $\le$ curvature''; each dies
for a distinct reason, and none of the reasons is repairable by sharpening the same tool: Brascamp--Lieb
is exact on the relevant test function so it has zero slack, curvature-dimension needs a constant that
is negative, RC-BL has a vanishing essential infimum on the bad set, sums of squares require a negative
variance, Cauchy--Schwarz is off by a factor, and the spectral route is an identity in disguise. Rows
7--12 close the structural variants: the two-halves split loses $E[\Phi'']$ exactly, the $\beta_{0}$
bridge fails its own acid test, pure-$\Phi$ quantities are excluded by the free-field test, $a$-uniform
statements die on a growing negative region, the fluctuation form is the target renamed, and the affine
form is an $O(a)$ approximation. What they jointly establish is negative and useful: the missing object
is not a better constant but a \emph{directional, conditioned, sign-carrying} lower bound on
$E[H_{vv}]$ at $a\neq0$.} \LN{}
\end{claimbox}

\src{Rows 1--10 are \code{LOOP-RUNBOOK.md} \S11.158 (the ten-row table, with the same numbers),
\S11.136 and \S11.150 (BL and the common root $E[H_{vv}]|_{a\neq0}<0$), \S11.149 (SOS) and \S11.152
(the corrected constant $C=3.95$ that supersedes both $1.7454$ and $10.05$); rows 11--12 are this
day's additions, \code{R09-FLUCT.md} \S2--\S4 and \code{R07-AFFINE.md} \S3--\S6; row 8's
counterexample is \code{P20-THETA.md} \S6 item 5, which also records that the same arithmetic
question was caught independently in \code{Q16-CONFLICT.md}.}

\subsection{Where strictness stops}
\label{sec:h1s-limit}

\paragraph{The volume.} $L=2$ has $192=4L^{4}\cdot3$ real link components; quotienting by the gauge
group leaves $144$; conditioning on the coarse mode $\langle v,X\rangle=a_{0}$ is a single scalar
constraint, so the fibre is
\begin{equation}
\dim=192-48-1=143 ,
\label{eq:h1s-143}
\end{equation}
and $143\gg20$ is the criterion for ``no interval arithmetic''. Symmetry does not help, and this is a
theorem rather than a disappointment: $W_{4}$ (order $384$) and the space group (order $6144$) are
\emph{finite}, so every orbit is a finite set and $\dim(M/G)=\dim M$; the only thing a finite group buys
is the measure factor $1/6144$ and a $3$-dimensional invariant tangent space at the flat point --- which
is the free field, where $G''\equiv0$ and the acid test fails. The per-link $1$-dimensional
factorization that made \S\ref{sec:h1s-cert} possible does not exist here, because $S_{W}$ is a sum over
plaquettes and the $\Phi$ terms share the links; the exact integration-by-parts pairing that gives
$NS=1/4$ is an accident of one plaquette. What $L=2$ has instead is a statistical certificate whose
price is explicit: a bounded-part-plus-tail-volume Hoeffding bound needs $N\approx3\times10^{8}$
independent samples ($\approx35$ days on one core), so the honest grade is \S\ref{sec:h1s-lattice}'s
\LE. One precise simplification does exist, at the corner $a_{0}=16\pi$: there $\rho_{v}\propto
\varepsilon^{3\cdot64-1}=\varepsilon^{191}$ and $G''=+191/\varepsilon^{2}\to+\infty$, which is
$Q08$'s $+3/(2\pi-a)^{2}$ promoted to $L=2$ --- but the corner is $16\pi$ away from the conditioning
point and says nothing about $a_{0}=0$. \LN

\src{\code{R04-L2CERT.md} \S0--\S2 (the dimension ledger, the rank checks, the finite-group theorem,
the orbit counts $8/4/1$ and the $12$-dimensional toron zero-mode decomposition
$192=45+12+135$), \S3 (the cost of a rigorous statistical certificate), \S4--\S4.1 (that $\theta=\pi$
does not simplify at $L=2$, and the corner law $3n-1$ with $n=64$);
\code{O24-L2-PUSH.md}:49 for the volume bookkeeping; \code{Q08-LARGEA.md} \S2.1 for the
single-plaquette version.}

\paragraph{The group.} $\mathrm{SU}(3)$ is measured and not proved, and the two numbers must be quoted
together because their signs differ:
\begin{equation}
\mathrm{gap}\big|_{\text{weight }|\Delta|^{2}e^{-S_W}}=+0.10172408,\qquad
\mathrm{gap}\big|_{\text{pure Haar}}=-0.23580249
\qquad(\theta=\pi,\ \beta=2.5,\ a=0) .
\label{eq:h1s-su3}
\end{equation}
The two earlier attempts are both dead, and for reasons worth keeping. The first concluded $S\equiv D$
from $S_{W}\equiv\beta$ on the slice; that inference is invalid, because $S_{W}\equiv\beta$ is a
\emph{value} and $G=\partial_{v}S_{W}$ is a derivative --- verified: $S_{W}-\beta=0.0$ at three points
while $G=\beta sg(u)=1.22167/0.97087/0.24784$, matching finite differences to $18$ digits. The second
error in the same file is the choice of $\Phi$: it used $-\log|\Delta|^{2}$, whereas the corpus object
is $\Phi=\sum_{i<j}g(d_{ij})$; the two differ by $\psi=\sum_{i<j}\log d_{ij}^{2}$, and with the wrong
choice the identity $E[\Phi'^{2}]=E[\Phi'']$ holds exactly ($48=48$ in exact rationals, interval width
$1.2\times10^{-119}$) and forces $S-D=-\operatorname{Var}(G)-2\operatorname{Cov}\le0$. The obstacle to
a certificate is dimension, not precision: the $a=0$ slice is $7$-dimensional ($1$ eigenvalue
direction times a $5$-dimensional flag), $2$D Weyl integration is exact only for class functions while
$G$ and $\Phi'$ are not, and $\Phi$ contains $\log$ and $\cot$/$\csc^{2}$ terms that destroy the
trigonometric-polynomial route that worked at $N=2$. \LE{}/~\LM

\src{\code{R12-SU3CORR.md} \S0--\S3 (both numbers, the $G\neq0$ verification, the two errors of the
earlier file, the exact $48=48$, and the $7$-dimensional obstacle), \S4 (the clause-by-clause verdict);
\code{R11-SU3ADJ.md} (the adjudication that the $S\equiv D$ file's inference is invalid);
\code{R08-SU3CERT.md} (the file under adjudication, kept on the record).}

\paragraph{The amplitude.} Lemma 8 is the whole of what is known off $a=0$: an exact splitting with an
explicit odd remainder. There is no $a$-uniform criterion, and row 10 of \S\ref{sec:h1s-roads} says why
one should not be expected from these legs: the negative region of $G''$ in $a$ widens and deepens
monotonically with $\beta$.

\paragraph{Two bookkeeping corrections made today, recorded because both changed a quoted number.}
(i) The $\theta=0$ certificate value was $1.04\times10^{-9}$ too low with an enclosure that did not
contain the truth; the corrected value is \eqref{eq:h1s-points}. (ii) The dimension bookkeeping that
called $L=2$ a $16$-link, $2$-plaquette lattice is wrong: $4L^{4}=64$ links and $6L^{4}=96$
plaquettes, and the $16$-link object is the blocked lattice $L_{c}=1$ obtained from $L=2$ by a factor
$2$. Neither correction moves a sign, and both are of the same family as the twenty-four rows of the
errata of the campaign: \emph{a quantity placed where it does not belong}.

\src{\code{R10-AUDITCERT.md} \S0 and \S3 for (i); \code{R04-L2CERT.md} \S1.0 for (ii), with
\code{Q11-CERT.md} \S6 (line 206) as the text corrected, and \code{O24-L2-PUSH.md}:49 as the
independent confirmation of $16$ sites / $64$ links / $96$ plaquettes.}

\paragraph{One further correction, which moves no sign.} (iii) $E[\kappa^{2}]=\infty$: the second moment
of $\kappa=\partial_{v}^{2}V$ diverges, so the ratio $\Delta/\sigma$ --- the $D/S$ quoted in
\S\ref{sec:h1s-lattice} --- has no numerical value that means anything, and its finite-difference quotes
carry an infinite-variance error bar in place of a finite one. The \emph{sign} conclusion stands: no sign
in this chapter is read off $\Delta/\sigma$, every one of them is a sign of $G''(0)$ itself. Same family
as (i)--(ii): \emph{a quantity placed where it does not belong}.

\subsection{Verdict of this chapter}
\label{sec:h1s-verdict}

\begin{itemize}
\item \LT{} \emph{One certificate, three values.} On a single plaquette, $\mathrm{SU}(2)$, $\beta=2.5$,
$a=0$: $\mathrm{gap}(0)=0.729924128620548449$, $\mathrm{gap}(\pi/2)=0.5429439223536897638$,
$\mathrm{gap}(\pi)=0.1359947110419882110$, each enclosed within $\pm10^{-55}$ by exact-rational
Taylor coefficients, directed rounding and Cauchy tails \eqref{eq:h1s-points}. This is the first
non-trivial positivity of $G''(0)$ in the campaign that is neither a Monte Carlo nor a
``high-precision estimates''.
\item \LC{} \emph{The whole $\theta$ interval, conditionally.} The identities
\eqref{eq:h1s-expfam}--\eqref{eq:h1s-mono} are proved and give four of the five terms of
$\partial_{c}\mathrm{gap}$ the right sign; the fifth is open, and on $501$ points the total is positive
with a minimum $0.4436$. The complement is cheap and explicit: four more certified points at rational
$c$ plus one bound on $|\mathrm{gap}''|$ cover $[0,\pi]$ unconditionally.
\item \LT{} \emph{Eight lemmas.} $\Phi''\succeq(N/12)\mathrm{id}$ for all $X$ and all directions,
sharp at $X=0$; $E[H_{vv}]\ge0$ for $\theta\le\pi$; $E[H_{vv}]=E[G_{v}^{2}]$ (Stein) with the spectral
route thereby empty; $H=\tfrac12\cot(\theta/2)G$ pointwise; $\operatorname{Var}(w)\le\tfrac12\langle w\rangle$
along three independent paths; $\langle w\rangle-2\operatorname{Var}\ge2\langle w\rangle^{2}$; the
skeleton identity \eqref{eq:h1s-skeleton}; and the exact splitting
$H=K_{0}G+\beta\Delta$.
\item \LT{}/\LC{} \emph{One window at two grades.} $\beta\sin^{2}(\theta/2)\le\cos(\theta/2)$ is a
theorem and covers $\theta\le0.388\pi$ at $\beta=2.5$ and $0.256\pi$ at $\beta=6$; the measured
$R_{\min}=32.3177$ widens this to $0.904\pi$ and $0.786\pi$ and the widening is one unproved
monotonicity away from being a theorem. The window is provably not $\theta$-uniform
($\beta_{0}\asymp8/(9(\pi-\theta)^{3})$) and at $\theta=\pi$ its premises collapse --- which is the
slice where the exact counterexample lives, at $\beta>4.30661$, i.e.\ at a different coupling, not a
different point.
\item \LE{}/\LM{} \emph{The lattice.} $G''(0)>0$ at every measured coupling: $+0.2835$, $+0.4116$,
$+0.5108$, $+0.6018$ at $L=2$ with $16/16$ and $8/8$ chains, and $+0.3396$ at $L=3$ with
$5.80\sigma$; no volume dependence at the $0.85\sigma$/$0.7\%$ level. At $L=2$, $\beta=6$ the five
estimators span $4.1$--$10.2\sigma$, and the only one below $5\sigma$ is a formula mismatch whose
correction returns $7.3\sigma$. The rank of the evidence is exactly \LE: one window, one mode, no
$a\to0$.
\item \LN{} \emph{Twelve roads.} Each ends on a number: $38$--$44\times$; a tautology at
$2\times10^{-16}$; $\mathrm{ess\,inf}K_{+}=0$; a required negative variance $-0.0367$; $2.6\times$
against $40$--$125\times$; a spectral identity; an exactly zero margin; a broken small-$B$ bridge at
$0.082919>0.0743$; a $65\%$ free-field remnant; a negative region growing $1.09\to5.66$; an exact
failure $0.027648>0$; and a residual $O(a)$ at $1.43\times10^{-3}$.
\item \LM{}/\LN{} \emph{What is left.} A directional, conditioned, sign-carrying lower bound on
$E[H_{vv}]$ at $a\neq0$; a strict certificate at $L=2$, which is $143$-dimensional; and any strict
statement at $\mathrm{SU}(3)$, which is $7$-dimensional on this slice. None of the three is buyable
with precision or with statistics.
\end{itemize}

\textbf{Not done, and not to be read as done.} No number in this chapter carries an $a\to0$ scaling;
the sharp variance leg ($\Psi\le1$) is measured on a grid, not proved, so the window's width at large
$\theta$ is a statement about a measured constant; the monotonicity that promotes three points to an
interval is proved in four of its five terms and measured in the fifth; $E[H_{vv}]\ge0$ is true only on
$\theta\le\pi$; the floor $\Phi''\succeq N/12$ cannot cancel the Wilson-side gap at $a\neq0$; the
$L=2$ and $L=3$ numbers are Monte Carlo on one mode with a finite-difference definition of $\kappa$;
and the $\mathrm{SU}(3)$ gap changes sign with the weight, so it is not even a well-defined number until
the reduction is fixed. \textbf{Clay remains $0\%$: this chapter states where $H1$ stands, and where it
stands is --- one certificate at three points with an interval conditional on one monotonicity, eight
strict lemmas with their scopes attached, one window with two boundaries and a third that is false, one
lattice ledger, twelve roads that end, and three walls named by their dimension.}
